% Pursuing Stacks (cotes 134-2 to 134-8): archive cote 134-7, site page 94 (PDF page 95), typescript page 536.
% typescript. Draft compared: ps6.tex:4724-4796.
% AI-transcribed by Claude Opus 5.5 (claude-opus-5-5, adaptive thinking, effort high), Message Batches API, batch ps4 (msgbatch_012Aoc1NuskATPsyBeKSYK7J), 2026-11-12; not yet checked by eye.
\scanpage{134-7/94}
\mspage{536}

encompassing choice for such objects seems to be sheaves on the category
of all \xout{\unsure{pre}}schemes  over k (or equivalently,  of all affine schemes over
k),  for a suitable topology such as the fpqc topology  (compare section
111,  pages 446-447\hand{)}\unsure{r}.  Apart from the technico-logical nuisance of this
not being a $\underline{U}$-site,  where $\underline{U}$ is our basic universe  (cf. p. 492\hand{)}\unsure{r},
which we'll ignore here (as it is'nt really too serious a difficulty),
\inserted{when} working with semisimplicial sheaves,  the embryo of foundational work
of Illusie's on the derived category of a topos via semisimplicial
sheaves\unsure{x}\hand{[} \xout{\unsure{x}} chapter 1\hand{]} becomes available.  Thus for a map of semi-\margin{x}
simplicial sheaves over k ,
\[
\text{(1)}\qquad X_{*} \longrightarrow Y_{*} \quad ,
\]
we know already what it means to be a quasi-isomorphism\inserted{(= weak equivalence)}. It seems like-
ly that,  even when drastically restricting the sheaves allowed as com-
ponents for our semisimplicial models,  the notion of quasi-isomorphism
relevant for these models should be the same as Illusie's\xout{\unsure{,}}  for the
encompassing topos.  This is one point we have been neglecting so far.
I am not going here to recall Illusie's definition (in terms of the
homotopy sheaves  $\underline{\pi}_{i}(X_{*}\ , s)$  where s is a section of $X_{*}$ ),  or trans-
late it into cohomological terms.   One point I want to make here,  is
that it is by no means automatic  that \xout{the} if (1) is a \xout{quasi-isomorphism}
weak equivalence,  the same holds for the corresponding map of semisim-
plicial sets
\[
\text{(2)}\qquad X_{*}(k) \longrightarrow Y_{*}(k) \quad .
\]
Thus \xout{\unsure{t}}there is no \xout{\unsure{tau}} more or less tautological ``sections'' functor from
Illusie's derived category\xout{\unsure{,}} to the category (Hot) of usual homotopy
types. \hand{(}There surely \uline{is} \xout{such} a \inserted[typed]{canonical} functor though,  in the case of a general
topos \inserted{T} and Illusie\unsure{\hand{'}}s corresponding d\xout{\unsure{r}}erived category,  namely what we would
like to call the ``\uline{cointegration over} T ''  \uline{functor},  which should come
out of the formalism of stacks we hav'nt \xout{yet} begun to  develop yet.\hand{)}
Presumably\inserted{however}\hand{,}  when working with \xout{models} semisimplicial models whose compo-
nents are restricted  to be unipotent bundles or something pretty close
to these (see below for examples),  whenever (1) is a weak equivalence,
the same \inserted{will} hold\deleted{s} for (2) .  As the choice of objects (let's call them
simply the ``bundles'') we are allowing,  should clearly be stable under
ring extension,  it will then follow that more generally,  for any
algebra k' over k,  the corresponding map
\[
\text{(2')}\qquad X_{*}(k') \longrightarrow Y_{*}(k')
\]

\margin{94}
